\documentclass{article}
\usepackage{amsmath}
\usepackage{graphicx}
\graphicspath{ {images/} }
\usepackage[spanish]{babel}
\usepackage{bbm}
\usepackage{amsthm}
\usepackage{authblk}
\usepackage{hyperref}
\usepackage[utf8]{inputenc}
\usepackage{mathrsfs}
\usepackage{amsfonts}
\usepackage{amssymb}
\usepackage{enumerate}
\usepackage[left=3cm,right=2cm,top=2cm,bottom=2cm]{geometry}
\usepackage{epsfig}
 \renewcommand{\baselinestretch}{0.93}
 \thispagestyle{empty}
 \pagestyle{empty}
\setlength{\textheight}{53pc} \setlength{\textwidth} {36pc}
\setlength{\topmargin}{-4mm}
\usepackage{multicol}
\newcommand*{\email}[1]{%
    \normalsize\href{mailto:#1}{#1}\par
    }
\setlength{\columnsep}{1cm}
\newcommand{\N}{\mathbb{N}}
\newcommand{\calM}{\cal{M}}
\newcommand{\calP}{\cal{P}}
\newcommand{\F}{\mathbb{F}}
\newcommand{\R}{\mathbb{R}}
\newcommand{\Z}{\mathbb{Z}}
\newcommand{\Q}{\mathbb{Q}}
\newcommand{\C}{\mathbb{C}}
\newcommand{\bbH}{\mathbb{H}}


\theoremstyle{definition}
\newtheorem{ejer}{Ejercicio}
\author{Marcos Morales}
\affil{Universidad Finis Terrae}
\title{Primera Semana\\}



	


\begin{document}


\begin{picture}(400, 75)
   \put(-15,15){\includegraphics[width=5cm]{UFT.png}}
\end{picture}



\begin{center}
\textbf{\Large{Taller 1}} \\
\end{center}

\setlength{\parindent}{0cm}



\begin{center}
\textbf{Cálculo Vectorial}
\end{center}

\hspace{100cm} \\ 


\textbf{Ejercicio 1:} Considere los rectas cuyas ecuaciones vectoriales son

\begin{align*}
 L_1 &: (1,2,0) + t(2,-1,3), \hspace{0.2cm}\text{$ t \in \mathbb{R}$}\\
 L_2 &:  (3,1,3) +\lambda (-1,2,-1), \hspace{0.2cm} \text{$ \lambda \in \mathbb{R}$} \\
 L_3 &:  (5,0,6)+ s(2,-4,2), \hspace{0.2cm}  \text{$ s \in \mathbb{R}$}\\
\end{align*}

a)  Encuentre las ecuaciones generales de cada recta.\\ 

b) Demuestre que las rectas $L_1$ y $L_2$ se intersectan y encuentre el punto de intersección. \\ 

c) Demuestre que las rectas $L_1$ y $L_3$ se intersectan y encuentre el punto de intersección.\\ 

d) Demuestre que las rectas $L_2$ y $L_3$  no se intersectan.\\ 

e) Determine el ángulo entre las rectas $L_1$ y $L_2$ y también el ángulo entre las rectas  $L_1$ y $L_3$. \\ \\ 


\textbf{Respuesta: } a) Las ecuaciones generales vienen dadas por \\

Para $L_1$

\begin{align*}
    \frac{x-1}{2} = \frac{y-2}{-1} = \frac{z}{3}
\end{align*}

Para $L_2$

\begin{align*}
    \frac{x-3}{-1} = \frac{y-1}{2} = \frac{z-3}{-1}
\end{align*}

Para $L_3$

\begin{align*}
    \frac{x-5}{2} = \frac{y}{-4} = \frac{z-6}{2}
\end{align*}


b) Tenemos que $\displaystyle \frac{d_1}{v_1} = \frac{2}{-1}\neq \frac{-1}{2} = \frac{d_2}{v_2}$. Creamos la matriz y calculamos el determinante 

\begin{center}
$\displaystyle \begin{vmatrix}
p_1-u_1 & d_1 & v_1 \\
p_2-u_2 & d_2 & v_2 \\
p_3-u_3 & d_3 & v_3 \\
\end{vmatrix} = \displaystyle \begin{vmatrix}
2 & 2 & -1 \\
-1 & -1 & 2 \\
3 & 3 & -1 \\
\end{vmatrix} = 0$
\end{center}


Por lo tanto las rectas se intersectan. Para calcular la intersección consideramos las ecuaciones paramétricas \\

Para $L_1$

\begin{align*}
    x &= 2t+1 \\
    y &= 2-t \\
    z &= 3t \\
\end{align*}

Para $L_2$

\begin{align*}
    x &= 3-\lambda \\
    y &= 2\lambda+1 \\
    z &= 3-\lambda \\
\end{align*}

Igualamos coordenada a coordenada (solo en $x$ y en $y$) y obtenemos


\begin{align*}
    2t+1 &= 3-\lambda \\
    2-t &= 2\lambda+1
\end{align*}

multiplicando la segunda por 2 obtenemos

\begin{align*}
    2t+1 &= 3-\lambda \\
    4-2t &= 4\lambda+2
\end{align*}

Sumando

\begin{align*}
5 &= 5+3\lambda
\end{align*}

Con lo cual se conluye que $\lambda=0$, luego el punto de intersección es $(3,1,3)$. \\ \\

c) Tenemos que $\displaystyle \frac{d_1}{v_1} = \frac{2}{2}\neq \frac{3}{2}$. Creamos la matriz y calculamos el determinante 

\begin{center}
$\displaystyle \begin{vmatrix}
p_1-u_1 & d_1 & v_1 \\
p_2-u_2 & d_2 & v_2 \\
p_3-u_3 & d_3 & v_3 \\
\end{vmatrix} = \displaystyle \begin{vmatrix}
4 & 2 & 2 \\
-2 & -1 & -4 \\
6 & 3 & 2 \\
\end{vmatrix} = 0$
\end{center}


Por lo tanto las rectas se intersectan. Para calcular la intersección consideramos las ecuaciones paramétricas, la de $L_1$ ya la tenemos \\

Para $L_3$

\begin{align*}
    x &= 2s+5 \\
    y &= -4s \\
    z &= 2s+6 \\
\end{align*}

Igualamos coordenada a coordenada (solo en $x$ y en $y$) y obtenemos


\begin{align*}
    2t+1 &= 2s+5 \\
    2-t &= -4s
\end{align*}

multiplicando la segunda por 2 obtenemos

\begin{align*}
    2t+1 &= 2s+5 \\
    4-2t &= -8s
\end{align*}

Sumando

\begin{align*}
5 &= 5-6s
\end{align*}

Con lo cual se conluye que $s=0$, luego el punto de intersección es $(5,0,6)$. \\ \\


d) Tenemos que $\displaystyle \frac{d_1}{v_1} =\frac{d_2}{v_2}= \frac{d_3}{v_3} = -\frac{1}{2}$. Por otra parte

\begin{align*}
    \frac{3-5}{-1}=2 \neq \frac{1-0}{2} =\frac{1}{2}
\end{align*}

Por lo tanto las rectas son paralelas y no se intersectan \\

e) El ángulo entre $L_1$ y $L_2$ viene dado por

\begin{align*}
    \alpha &= \arccos\left( \frac{(2,-1,3) \cdot (-1,2,-1)}{||(2,-1,3)|| || (-1,2,-1)||}\right) \\
    &= \arccos\left( -\frac{7}{\sqrt{14}\sqrt{6}}\right) \\
\end{align*}


El ángulo entre $L_1$ y $L_2$ viene dado por

\begin{align*}
    \beta &= \arccos\left( \frac{(2,-1,3) \cdot (2,-4,2)}{||(2,-1,3)|| || (2,-4,2)||}\right) \\
    &= \arccos\left( \frac{14}{\sqrt{14}\sqrt{24}}\right) \\
\end{align*}



\newpage


\textbf{Ejercicio 2:} Considere los planos dados por las ecuaciones

\begin{align*}
    \Pi_1 &: -3x+5y+2z-3k = 0 \\
    \Pi_2 &: kx+ y-7z +4 =0
\end{align*}

Donde $k$ es una constante. \\

a)  Determine el valor de la constante $k$ para que los planos sean perpendiculares. \\ 

b) Determine la intersección entre los dos planos. \\ \\ 


\textbf{Solución:} a) Los vectores normales son

\begin{align*}
    n_1 &= (-3,5,2) \\
    n_2 &= (k,1,-7)
\end{align*}

Al hacer el producto punto, obtenemos

\begin{align*}
    n_1 \cdot n_2 = -3k+5-14
\end{align*}

Como deben ser perpendiculares se debe cumplir que

\begin{align*}
    -3k+5-14 &=0 \\
    k &= -3
\end{align*}

b) Los planos son

\begin{align*}
    \Pi_1 &: -3x+5y+2z+9 = 0 \\
    \Pi_2 &: -3x+ y-7z +4 =0
\end{align*}

Despejamos $y$ en la segunda ecuación, obtenemos

\begin{align*}
    y =3x+7z-4
\end{align*}

Reeemplazando en la primera ecuación

\begin{align*}
    -3x+5(3x+7z-4)+2z+9 = 0 
\end{align*}

Ahora despejamos $x$, obtenemos


\begin{align*}
    -3x+15x+35z-20+2z+9 &= 0 \\
    x = \frac{11-37z}{12}
\end{align*}

Reemplzando en $y$, obtenemos 

\begin{align*}
    y &=3\left( \frac{11-37z}{12} \right)+7z-4 \\
    &= \frac{-5-9z}{4} \\
\end{align*}

Luego, haciendo $z=t$, la parametrización es la siguiente


\begin{align*}
    \left(  \frac{11-37t}{12}, \frac{-5-9t}{4}, t\right) 
\end{align*}

O bien

\begin{align*}
    \left(  \frac{11}{12}, \frac{-5}{4}, 0\right) +t \left( -\frac{37}{12}, -\frac{9}{4},1\right) 
\end{align*}

que es la ecuación vectorial de una recta.











\newpage

\textbf{Ejercicio 3:} Considere el campo escalar

\begin{align*}
    f(x,y) = \ln(x^2+y^2-4)
\end{align*}

a) Determine el dominio de $f$. \\

b) Halle las curvas de nivel  $C_0(f)$, $C_1(f)$ y $C_2(f)$. \\

c) Dibuje el dominio de $f$ y las curvas de nivel de la parte $(b)$. \\


\textbf{Solución:} a) Tenemos que

\begin{align*}
    \text{Dom}(f) &= \lbrace (x,y) \in \mathbb{R}^2 : x^2+y^2-4 > 0 \rbrace \\
    &= \lbrace (x,y) \in \mathbb{R}^2 : x^2+y^2> 4 \rbrace \\
\end{align*}


b) en general tenemos para un valor $c$

\begin{align*}
    c= \ln(x^2+y^2-4)
\end{align*}

Luego

\begin{align*}
    e^c= x^2+y^2-4
\end{align*}

Luego

\begin{align*}
     x^2+y^2 = 4+e^c
\end{align*}

\begin{align*}
     x^2+y^2 = (\sqrt{4+e^c})^2
\end{align*}

Para $C=0$, obtenemos

\begin{align*}
     x^2+y^2 = (\sqrt{5})^2
\end{align*}


Para $C=1$, obtenemos

\begin{align*}
     x^2+y^2 = (\sqrt{4+e})^2
\end{align*}

Para $C=2$, obtenemos


\begin{align*}
     x^2+y^2 = (\sqrt{4+e^2})^2
\end{align*}


\newpage

c) El dominio correspondde a la parte extrena del círculo de radio 4 \\

Cada curva de nivel son círculos, de radio $\sqrt{5}$, $\sqrt{4+e}$ y $\sqrt{4+e^2}$




\end{document} 